% appD-source-index.tex — 附录 D 源码索引 v1（RFC-A.5）
% permalink 基准 commit：6cd32de（tag v4.0.0，2026-09-17；anchors.yaml repo_commit）
\chapter{源码索引}\label{app:D}

正文关键代码段的完整实现入口。行区间与 \texttt{scripts/\allowbreak{}anchors.\allowbreak{}yaml} 冻结登记一致；
链接为 GitHub commit permalink（防 branch 漂移）。

\begin{longtable}{llp{7.2cm}}
\toprule
章 & 主题 & 源码区间（permalink） \\
\midrule
\endfirsthead
\toprule
章 & 主题 & 源码区间（permalink） \\
\midrule
\endhead
A & common.cuh 工具箱 & \href{https://github.com/xlite-dev/LeetCUDA/blob/6cd32deca683d6c4aae2011988c9f4a35bc5f528/kernels/interview/common.cuh#L1-L803}{common.cuh L1--L803} \\
第01 & GPU 架构/Roofline 速查 & \href{https://github.com/xlite-dev/LeetCUDA/blob/6cd32deca683d6c4aae2011988c9f4a35bc5f528/kernels/interview/base.cuh#L1-L86}{base.cuh L1--L86} \\
第02 & warp/block 归约与 dot & \href{https://github.com/xlite-dev/LeetCUDA/blob/6cd32deca683d6c4aae2011988c9f4a35bc5f528/kernels/interview/base.cuh#L87-L305}{base.cuh L87--L305} \\
第03 & 向量化与原子操作 & \href{https://github.com/xlite-dev/LeetCUDA/blob/6cd32deca683d6c4aae2011988c9f4a35bc5f528/kernels/interview/base.cuh#L306-L386}{base.cuh L306--L386} \\
第04 & softmax 三级递进 & \href{https://github.com/xlite-dev/LeetCUDA/blob/6cd32deca683d6c4aae2011988c9f4a35bc5f528/kernels/interview/base.cuh#L520-L665}{base.cuh L520--L665} \\
第05 & merge\_attn\_states & \href{https://github.com/xlite-dev/LeetCUDA/blob/6cd32deca683d6c4aae2011988c9f4a35bc5f528/kernels/interview/base.cuh#L387-L519}{base.cuh L387--L519} \\
第06 & RMSNorm/LayerNorm & \href{https://github.com/xlite-dev/LeetCUDA/blob/6cd32deca683d6c4aae2011988c9f4a35bc5f528/kernels/interview/base.cuh#L666-L801}{base.cuh L666--L801} \\
第07 & RoPE 与转置 & \href{https://github.com/xlite-dev/LeetCUDA/blob/6cd32deca683d6c4aae2011988c9f4a35bc5f528/kernels/interview/base.cuh#L802-L909}{base.cuh L802--L909} \\
第08 & sgemv 三种划分 & \href{https://github.com/xlite-dev/LeetCUDA/blob/6cd32deca683d6c4aae2011988c9f4a35bc5f528/kernels/interview/sgemv.cuh#L1-L102}{sgemv.cuh L1--L102} \\
第09 & sgemm 阶梯一 & \href{https://github.com/xlite-dev/LeetCUDA/blob/6cd32deca683d6c4aae2011988c9f4a35bc5f528/kernels/interview/sgemm.cuh#L6-L183}{sgemm.cuh L6--L183} \\
第10 & sgemm TF32/WMMA & \href{https://github.com/xlite-dev/LeetCUDA/blob/6cd32deca683d6c4aae2011988c9f4a35bc5f528/kernels/interview/sgemm.cuh#L184-L434}{sgemm.cuh L184--L434} \\
第11 & hgemm mma.sync & \href{https://github.com/xlite-dev/LeetCUDA/blob/6cd32deca683d6c4aae2011988c9f4a35bc5f528/kernels/interview/hgemm.cuh#L3-L397}{hgemm.cuh L3--L397} \\
第12 & hgemm swizzle & \href{https://github.com/xlite-dev/LeetCUDA/blob/6cd32deca683d6c4aae2011988c9f4a35bc5f528/kernels/interview/hgemm.cuh#L399-L716}{hgemm.cuh L399--L716}, \href{https://github.com/xlite-dev/LeetCUDA/blob/6cd32deca683d6c4aae2011988c9f4a35bc5f528/kernels/interview/common.cuh#L110-L348}{common.cuh L110--L348} \\
第13 & WGMMA/TMA/mbarrier & \href{https://github.com/xlite-dev/LeetCUDA/blob/6cd32deca683d6c4aae2011988c9f4a35bc5f528/kernels/interview/hgemm.cuh#L1428-L1857}{hgemm.cuh L1428--L1857}, \href{https://github.com/xlite-dev/LeetCUDA/blob/6cd32deca683d6c4aae2011988c9f4a35bc5f528/kernels/interview/common.cuh#L350-L803}{common.cuh L350--L803} \\
第14 & SM120 TMA+WS & \href{https://github.com/xlite-dev/LeetCUDA/blob/6cd32deca683d6c4aae2011988c9f4a35bc5f528/kernels/interview/hgemm.cuh#L1859-L2100}{hgemm.cuh L1859--L2100}, \href{https://github.com/xlite-dev/LeetCUDA/blob/6cd32deca683d6c4aae2011988c9f4a35bc5f528/kernels/interview/common.cuh#L462-L475}{common.cuh L462--L475} \\
第15 & FA 原理（头注释） & \href{https://github.com/xlite-dev/LeetCUDA/blob/6cd32deca683d6c4aae2011988c9f4a35bc5f528/kernels/interview/flash_attn.cuh#L5-L110}{flash\_attn.cuh L5--L110} \\
第16 & FA2 split-Q+MMA & \href{https://github.com/xlite-dev/LeetCUDA/blob/6cd32deca683d6c4aae2011988c9f4a35bc5f528/kernels/interview/flash_attn.cuh#L5-L790}{flash\_attn.cuh L5--L790} \\
第17 & FA2 TMA+WS & \href{https://github.com/xlite-dev/LeetCUDA/blob/6cd32deca683d6c4aae2011988c9f4a35bc5f528/kernels/interview/flash_attn.cuh#L792-L1440}{flash\_attn.cuh L792--L1440} \\
第18 & FA3 双 consumer & \href{https://github.com/xlite-dev/LeetCUDA/blob/6cd32deca683d6c4aae2011988c9f4a35bc5f528/kernels/interview/flash_attn.cuh#L1441-L2192}{flash\_attn.cuh L1441--L2192} \\
第19 & FFPA split-D 全篇 & \href{https://github.com/xlite-dev/LeetCUDA/blob/6cd32deca683d6c4aae2011988c9f4a35bc5f528/kernels/interview/ffpa_attn.cuh#L1-L641}{ffpa\_attn.cuh L1--L641} \\
第20 & CuTe Layout 对照 & \href{https://github.com/xlite-dev/LeetCUDA/blob/6cd32deca683d6c4aae2011988c9f4a35bc5f528/kernels/interview/hgemm.cuh#L788-L820}{hgemm.cuh L788--L820} \\
第21 & CuTe TiledCopy 对照 & \href{https://github.com/xlite-dev/LeetCUDA/blob/6cd32deca683d6c4aae2011988c9f4a35bc5f528/kernels/interview/hgemm.cuh#L821-L1196}{hgemm.cuh L821--L1196} \\
第22 & CuTe TiledMMA 对照 & \href{https://github.com/xlite-dev/LeetCUDA/blob/6cd32deca683d6c4aae2011988c9f4a35bc5f528/kernels/interview/ffpa_attn.cuh#L31-L79}{ffpa\_attn.cuh L31--L79} \\
第23 & CuTe SW128 atom 对照 & \href{https://github.com/xlite-dev/LeetCUDA/blob/6cd32deca683d6c4aae2011988c9f4a35bc5f528/kernels/interview/hgemm.cuh#L1232-L1270}{hgemm.cuh L1232--L1270} \\
第24 & CuTe HGEMM & \href{https://github.com/xlite-dev/LeetCUDA/blob/6cd32deca683d6c4aae2011988c9f4a35bc5f528/kernels/interview/hgemm.cuh#L718-L1427}{hgemm.cuh L718--L1427} \\
第25 & CuTe FA 三实现 & \href{https://github.com/xlite-dev/LeetCUDA/blob/6cd32deca683d6c4aae2011988c9f4a35bc5f528/kernels/interview/flash_attn.cuh#L2196-L3488}{flash\_attn.cuh L2196--L3488} \\
第26 & CuTe FFPA 双 TiledMMA & \href{https://github.com/xlite-dev/LeetCUDA/blob/6cd32deca683d6c4aae2011988c9f4a35bc5f528/kernels/interview/ffpa_attn.cuh#L30-L65}{ffpa\_attn.cuh L30--L65}, \href{https://github.com/xlite-dev/LeetCUDA/blob/6cd32deca683d6c4aae2011988c9f4a35bc5f528/kernels/interview/ffpa_attn.cuh#L79-L375}{ffpa\_attn.cuh L79--L375}, \href{https://github.com/xlite-dev/LeetCUDA/blob/6cd32deca683d6c4aae2011988c9f4a35bc5f528/kernels/interview/ffpa_attn.cuh#L377-L640}{ffpa\_attn.cuh L377--L640} \\
\multicolumn{3}{l}{\footnotesize 增补章（ch26b/ch26c/ch34/ch35）源码为本仓新增/追加区间（ffpa\_attn.cuh L643--L1306、} \\
\multicolumn{3}{l}{\footnotesize fp8\_gemm.cuh L1--L958、notes-v2.cu 集成 harness），随增补 commit 在 anchors.yaml 冻结} \\
\multicolumn{3}{l}{\footnotesize 以下 Part VI（第27--33 章）源码在 \texttt{ffpa-attn} 仓（commit \texttt{861d75e}），路径前缀以 \texttt{csrc/cuffpa/} 为主（bench 与 CLI 见第33章）} \\
第27 & 量化注意力数学 & \href{https://github.com/xlite-dev/ffpa-attn/blob/861d75ee87824a14a9ebfe295d54c7dac3e3cabc/csrc/cuffpa/cute/fp8/fp8_pscale.cuh#L17-L37}{cute/fp8/fp8\_pscale.cuh L17--L37}, \href{https://github.com/xlite-dev/ffpa-attn/blob/861d75ee87824a14a9ebfe295d54c7dac3e3cabc/csrc/cuffpa/cute/fp8/fp8_pscale.cuh#L111-L138}{L111--L138} \\
第28 & fp8 量化前处理链 & \href{https://github.com/xlite-dev/ffpa-attn/blob/861d75ee87824a14a9ebfe295d54c7dac3e3cabc/csrc/cuffpa/cute/fp8/smooth_k.cuh#L12-L84}{cute/fp8/smooth\_k.cuh L12--L84}, \href{https://github.com/xlite-dev/ffpa-attn/blob/861d75ee87824a14a9ebfe295d54c7dac3e3cabc/csrc/cuffpa/cute/fp8/quantize_fp8.cuh#L29-L201}{cute/fp8/quantize\_fp8.cuh L29--L201}, \href{https://github.com/xlite-dev/ffpa-attn/blob/861d75ee87824a14a9ebfe295d54c7dac3e3cabc/csrc/cuffpa/cute/hadamard.cuh#L1-L92}{cute/hadamard.cuh L1--L92} \\
第29 & fp8 persist-D & \href{https://github.com/xlite-dev/ffpa-attn/blob/861d75ee87824a14a9ebfe295d54c7dac3e3cabc/csrc/cuffpa/cute/fp8/sm_120/persist_d.cuh#L37-L215}{cute/fp8/sm\_120/persist\_d.cuh L37--L215}, \href{https://github.com/xlite-dev/ffpa-attn/blob/861d75ee87824a14a9ebfe295d54c7dac3e3cabc/csrc/cuffpa/cute/fp8/sm_120/persist_d.cuh#L583-L1031}{L583--L1031}, \href{https://github.com/xlite-dev/ffpa-attn/blob/861d75ee87824a14a9ebfe295d54c7dac3e3cabc/csrc/cuffpa/cute/fp8/reg2reg_8b.cuh#L14-L141}{cute/fp8/reg2reg\_8b.cuh L14--L141} \\
第30 & fp8 split-D/M4N2 & \href{https://github.com/xlite-dev/ffpa-attn/blob/861d75ee87824a14a9ebfe295d54c7dac3e3cabc/csrc/cuffpa/cute/fp8/sm_120/split_d.cuh#L39-L245}{cute/fp8/sm\_120/split\_d.cuh L39--L245}, \href{https://github.com/xlite-dev/ffpa-attn/blob/861d75ee87824a14a9ebfe295d54c7dac3e3cabc/csrc/cuffpa/cute/fp8/sm_120/split_d.cuh#L446-L756}{L446--L756}, \href{https://github.com/xlite-dev/ffpa-attn/blob/861d75ee87824a14a9ebfe295d54c7dac3e3cabc/csrc/cuffpa/cute/fp8/sm_120/split_d_m4n2.cuh#L300-L970}{split\_d\_m4n2.cuh L300--L970} \\
第31 & fp4 量化链 & \href{https://github.com/xlite-dev/ffpa-attn/blob/861d75ee87824a14a9ebfe295d54c7dac3e3cabc/csrc/cuffpa/cute/fp4/quantize_fp4.cuh#L32-L320}{cute/fp4/quantize\_fp4.cuh L32--L320}, \href{https://github.com/xlite-dev/ffpa-attn/blob/861d75ee87824a14a9ebfe295d54c7dac3e3cabc/csrc/cuffpa/cute/fp4/quantize_fp4.cuh#L322-L496}{L322--L496}, \href{https://github.com/xlite-dev/ffpa-attn/blob/861d75ee87824a14a9ebfe295d54c7dac3e3cabc/csrc/cuffpa/cute/fp4/delta_s.cuh#L2-L161}{cute/fp4/delta\_s.cuh L2--L161} \\
第32 & fp4 persist-D & \href{https://github.com/xlite-dev/ffpa-attn/blob/861d75ee87824a14a9ebfe295d54c7dac3e3cabc/csrc/cuffpa/cute/fp4/sm_120/persist_d.cuh#L1-L120}{cute/fp4/sm\_120/persist\_d.cuh L1--L120}, \href{https://github.com/xlite-dev/ffpa-attn/blob/861d75ee87824a14a9ebfe295d54c7dac3e3cabc/csrc/cuffpa/cute/fp4/sm_120/persist_d.cuh#L722-L1066}{L722--L1066}, \href{https://github.com/xlite-dev/ffpa-attn/blob/861d75ee87824a14a9ebfe295d54c7dac3e3cabc/csrc/cuffpa/cute/fp4/fp4_pscale.cuh#L117-L346}{cute/fp4/fp4\_pscale.cuh L117--L346} \\
第33 & bench 与 CLI & \href{https://github.com/xlite-dev/ffpa-attn/blob/861d75ee87824a14a9ebfe295d54c7dac3e3cabc/bench/bench_fp8.py#L1-L140}{bench/bench\_fp8.py L1--L140}, \href{https://github.com/xlite-dev/ffpa-attn/blob/861d75ee87824a14a9ebfe295d54c7dac3e3cabc/src/ffpa_attn/cli/_bench.py#L347-L490}{src/ffpa\_attn/cli/\_bench.py L347--L490} \\
\bottomrule
\end{longtable}

基准 commit：LeetCUDA \texttt{6cd32de}（v4.0.0，2026-09-17 冻结）；
Part VI 各章 ffpa-attn \texttt{861d75e}（2026-09-16 锚定）；Part V 两章
（第~\ref{ch:34}、\ref{ch:35}）与增补章（ch26b/ch26c）为本仓新增区间，
随增补 commit 在 \texttt{anchors.yaml} 冻结。行号漂移由
\texttt{verify\_\allowbreak{}anchors.\allowbreak{}py} 兜底：SHA256 + 首末行锚点双重校验。
